\section{The Problem Statement}

\subsection{Official description}

\begin{keybox}{SIH26117 --- Mangalore Refinery and Petrochemicals Limited}
\emph{Background.} Refineries, PSUs, defence-linked manufacturing units and
government offices generate a lot of routine but sensitive knowledge work:
approval notes, board presentations, engineering calculations, code for
internal tools, review of scanned drawings and inspection reports. None of this
can go through cloud AI assistants because the underlying data is confidential:
piping and instrumentation diagrams, financials, vendor negotiations,
unreleased designs, internal correspondence, confidential business strategies.
Company policy keeps this data on premises, so people either do the work
manually --- losing productivity --- or they quietly paste confidential
material into public tools anyway.
\end{keybox}

\subsection{Decomposition of the requirements}

The statement contains five explicit demonstration requirements. Each is
treated in this report as a verifiable acceptance criterion.

\begin{center}
\renewcommand{\arraystretch}{1.3}
\begin{tabularx}{\linewidth}{@{}c X l@{}}
\toprule
\textbf{\#} & \textbf{Requirement as stated} & \textbf{Section} \\
\midrule
1 & Model auto-selection across at least two different task types &
    \S\ref{sec:router} \\
2 & An agentic task carried end to end, e.g.\ reading a scanned inspection
    report, extracting findings, drafting an approval note as a Word file &
    \S\ref{sec:agent} \\
3 & A coding task run and verified in a sandbox & \S\ref{sec:sandbox} \\
4 & A multimodal task involving image or scanned document understanding &
    \S\ref{sec:multimodal} \\
5 & Evidence, through logs or a visible network monitor, that no external calls
    are made at any point & \S\ref{sec:sovereignty} \\
\bottomrule
\end{tabularx}
\end{center}

Alongside these, the statement imposes four architectural constraints:

\begin{enumerate}
\item \textbf{Not locked to one model.} Multiple open-weight models must be
      supported simultaneously, with automatic selection based on task needs.
\item \textbf{Extensible without redesign.} New open-weight models must be
      addable later, because the field moves quickly.
\item \textbf{Genuinely agentic.} The system must plan multi-step work, call
      local tools, and iterate --- not answer once and stop.
\item \textbf{Real deliverables.} Output must be approval notes, office
      documents, working code and calculations with steps shown --- not chat
      replies.
\end{enumerate}

\subsection{Hardware constraint}

The statement permits demonstration on ``a single workstation or server with a
mid-range GPU'', explicitly allowing a smaller open-weight model if 120B-class
hardware is unavailable at the venue. The reference implementation therefore
targets a single 16\,GB consumer card, and treats the ability to run on modest
hardware as a feature rather than a compromise: an organisation that must
procure an H100 before it can pilot the system will not pilot the system.

\subsection{Why this is hard}

\begin{description}[leftmargin=0pt,style=nextline]
\item[Memory is the binding constraint.] A 16\,GB card cannot hold a reasoning
model, a coding model and a vision model concurrently. Naive multi-model
serving thrashes or crashes.
\item[Proving a negative.] Demonstrating that \emph{no} external call occurred
is qualitatively harder than demonstrating a feature works. It requires
enforcement at a layer the application cannot bypass.
\item[Retrieval over structure.] Industrial questions are frequently about
dependency and impact. These are graph problems wearing the clothing of search
problems, and conventional retrieval-augmented generation answers them badly.
\item[Trust in output.] A system that fabricates a plausible approval note is
worse than no system. Every claim must carry provenance back to a source page.
\end{description}
